The subject of this article is to focus on axially symmetric perturbations of the Kerr metric. In precise terms, the Kerr metric $\bigl(\bar{M}, \bar{g}\bigr)$ can be represented in Boyer--Lindquist coordinates $(t, r, \theta, \phi)$ as
\begin{align}
\bar{g} &{}= - \biggl(\frac{\Delta - a^2 \sin^2 \theta}{\Sigma} \biggr) {\rm d}t^2 - \frac{2a \sin^2 \theta \bigl(r^2 + a^2 -\Delta\bigr)}{\Sigma} {\rm d}t {\rm d}\phi \notag\\
&\quad{}+ \biggl(\frac{\bigl(r^2 +a^2\bigr)^2 -\Delta a^2 \sin^2 \theta}{\Sigma}\biggr) \sin^2 \theta {\rm d}\phi^2 + \frac{\Sigma}{\Delta} {\rm d}r^2 + \Sigma {\rm d}\theta^2, \label{BL-Kerr}
\end{align}
where
\begin{align*}
&\Sigma := r^2 + a^2 \cos^2 \theta, \\
&\Delta := r^2 -2Mr + a^2 \quad \textnormal{with the real roots} \ \{r_-,r_+\}, \qquad
 r_{+} := M + \sqrt{M^2- a^2} > r_{-}
\end{align*}
and
\[
\theta \in [0, \pi],\qquad r \in (r_+, \infty),\qquad \phi \in [0, 2\pi).
\]
It is well known (cf.\ \cite{YCB-VM96, Ycb-Mon_01, diss_13}) that the Einstein equations \eqref{EVE} on spacetimes $\bigl(\bar{M}, \bar{g}\bigr)$ with one isometry $\bigl(\frac{\ptl}{\ptl \phi}\bigr)$, represented in Weyl--Papapetrou coordinates,
\begin{align*}
\bar{g} = {\rm e}^{-2\gamma} g + {\rm e}^{2\gamma} ({\rm d}\phi + \mathcal{A}_\nu {\rm d}x^\nu)^2, \qquad \nu= 0, 1, 2,
\end{align*}
where ${\rm e}^{2 \gamma}$ is the square of the norm of the rotational Killing vector $\ptl_\phi$, $\mathcal{A}$ is a $1$-form and $g$ is the induced metric on orbit space $M := \bar{M}/{\rm SO}(2)$, admit a dimensional reduction to a $(2 + 1)$-dimensional Einstein wave map system
\begin{subequations} \label{ewm-system}
\begin{align}
&E_{\mu \nu} = T_{\mu \nu}, \\
&\square_g U^A + \leftexp{(h)}{\Gamma}^A_{BC} g^{\mu \nu} \ptl_\mu U^B \ptl_\nu U^C= 0 \qquad \text{on} \ (M, g),
\end{align}
\end{subequations}
where $\square_g$ is the covariant wave operator, $E_{\mu \nu}$ is the Einstein tensor in the interior of the quotient $(M, g) := \bigl(\bar{M}, \bar{g}\bigr)/ {\rm SO}(2)$ and $T$ is the stress energy tensor of the wave map $U \colon (M, g) \to (\mathbb{N}, h)$, $\mathbb{N}$ is the negatively curved hyperbolic $2$-plane,
\begin{align*}
T_{\mu \nu} = \ip{\ptl_\mu U} {\ptl_\nu U }_{h(U)} - \halb g_{\mu \nu} \ip{\ptl_\sigma U}{ \ptl^\sigma U}_{h(U)}.
\end{align*}
Introducing the coordinates $(\r, z)$ such that $\r = R \sin \theta$, $z= R \cos
\theta$, where $R := \halb \bigl(r-M+ \sqrt{\Delta}\bigr)$, the Kerr metric \eqref{BL-Kerr} can be represented in the Weyl--Papapetrou form as
\begin{align*}
\bar{g} = \Sigma \zeta^{-1} \bigl(-\Delta {\rm d}t^2 + \zeta R^{-2} \bigl({\rm d}\r^2 + {\rm d}z^2\bigr)\bigr) + \sin^2\theta \Sigma^{-1} \zeta
\bigl({\rm d}\phi - 2aMr \zeta^{-1} {\rm d}t\bigr)\bigr)^2,
\end{align*}
where \smash{$\zeta = \bigl(r^2 +a^2\bigr)^2 - a^2 \Delta \sin^2 \theta$}. Furthermore, the Kerr metric can also be represented in the Weyl--Papapetrou form using functions $(\bar{\r}, \bar{z})$
such that
\begin{align*}
\bar{\r} = \r - \frac{\bigl(M^2 -a^2\bigr)}{4 R^2}\r, \qquad \text{and} \qquad \bar{z} = z + \frac{\bigl(M^2 -a^2\bigr)}{4 R^2}z
\end{align*}
(cf.\ \cite[Appendix A]{GM17_gentitle} for details). Now we shall turn to the axially symmetric perturbation theory of the Kerr metric.

In view of the peculiar behaviour of the $2 + 1$ Einstein-wave map system, a detailed discussion of our methods is relevant for our article and perhaps also interesting to the reader. Consider the Hamiltonian energy of an axially symmetric linear wave equation propagating on the Kerr metric ($\square_g u =0$),
\begin{align*}
H^{\text{LW}} := \int_{\olin{\Sigma}} \biggl(\halb \bar{N}\bar{\mu}^{-1}_q v^2 + vN^i \ptl_i u+ \halb N \bar{\mu}_q \bar{q}^{ij} \ptl_i u \ptl_j u \biggr) {\rm d}^3 x,
\end{align*}
where $v$ is the conjugate momentum of $u$, the energy is directly positive-definite.
However, this simplification does not carry forward to the Maxwell equations on the Kerr metric
\[
H^{{\rm Max}} := \int \biggl(\halb N \bar{q}_{ij} \bar{\mu}_q (\mathfrak{E}^i \mathfrak{E}^j + \mathfrak{B}^i \mathfrak{B}^j) + N^i \eps_{ijk}\mathfrak{E}^j \mathfrak{B}^k \biggr) {\rm d}^3 x,
\]
where
\[
\mathfrak{E}^i := \halb \eps^{ijk} \leftexp{*}F_{jk}, \qquad
\mathfrak{B}^i := \halb \eps^{ijk} F_{jk}
\]
are the electric and magnetic field densities respectively and $F$ is the Faraday tensor which satisfies the Maxwell equations
\begin{align*}
 {\rm d} \star F =0, \qquad {\rm d} F =0.
\end{align*}

Actually, one can construct counter examples of positivity of energy density, for instance, using time-symmetric Maxwell fields (cf.\ the discussion in \cite[Section 2]{GM17_gentitle}). In a crucial work, Dain--de Austria \cite{DA_14} had arrived at a positive-definite energy for the gravitational perturbations of extremal Kerr black holes using the Brill mass formula \cite{D09} and subsequent use of Carter's identity \cite{Car_71}, originally developed for black hole uniqueness theorems. In a Weyl coordinate system for the spacelike hypersurface $\bigl(\olin{\Sigma}, q\bigr)$ in extremal Kerr spacetime, their positive-definite energy for axially symmetric perturbations is obtained
from perturbations of the Brill mass formula, which in turn is obtained from multiplying a factor with the Hamiltonian constraint that conveniently results in a volume form (in the chosen Weyl coordinate system) that is useful in its representation.

In order to construct a positive-definite energy for the perturbations of Kerr--Newman metric for the full-subextremal range, we delve into the variational structure of the relevant field equations. The beautiful linearization stability framework, developed by V.~Moncrief, J.~Marsden and A.~Fischer~\cite{FM_75,FMM_80,Mon_75}, provides a
natural mechanism to construct an energy-functional based on the kernel of the adjoint of the deformations around the Kerr metric of the dimensionally reduced constraint map. This recognition allows us to extend results to the full sub-extremal range $(\vert a \vert, \vert Q \vert <M)$ of the perturbations of the Kerr--Newman metric \cite{GM17_gentitle}, which is a solution of Einstein--Maxwell equations of general relativity.

Let us briefly outline our construction of a positive-definite energy functional for axially symmetric perturbations of Kerr black hole spacetimes. Consider the ADM decomposition of $\bar{M}=\olin{\Sigma} \times \mathbb{R}$. Suppose the group ${\rm SO}(2)$ acts on $\olin{\Sigma}$ through isometries. Let $\Gamma$ be the non-empty fixed-point set. Suppose the norm squared of the Killing vector generating the rotational isometry is denoted by ${\rm e}^{2\gamma}$. In the dimensional reduction ansatz, the metric $\bar{g}$ is
\begin{align}\label{KK-ADM}
\bar{g} = {\rm e}^{-2\gamma} \bigl(-N^2 {\rm d}t^2 + q_{ab} ({\rm d}x^a + N^a {\rm d}t) \otimes \bigl({\rm d}x^b + N^b {\rm d}t\bigr)\bigr) + {\rm e}^{2\gamma} ({\rm d}\phi + \mathcal{A}_0 {\rm d}t + \mathcal{A}_a {\rm d}x^a)^2.
\end{align}
It may be noted that this metric form is combination of ADM formalism and Weyl--Papapetrou coordinates \cite{Moncrief_74}.
In the dimensional reduction framework, identifying the reduced conjugate momenta, which form the reduced phase space in $(M, g);$ and the corresponding reduced Hamiltonian formalism is nontrivial. This construction was done in \cite{kaluz1}. Define the conjugate momentum corresponding to the metric $q_{ab}$ as follows:
\begin{align*}
\mbo{\pi}^{ab} = {\rm e}^{-2\gamma} \bar{\mbo{\pi}}^{ab}, \qquad \bar{q}_{ab} = {\rm e}^{-2\gamma} q_{ab} + {\rm e}^{2\gamma} \mathcal{A}_a \mathcal{A}_b.
\end{align*}
As a consequence, the ADM action principle transforms to
\begin{align*}
J = \int^{t_2}_{t_1}\int_{\Sigma} \bigl(\mbo{\pi}^{ab} \ptl_t q_{ab} + \mathcal{E}^a \ptl_t \mathcal{A}_a + p \ptl_t \gamma - N H - N^aH_a+ \mathcal{A}_0 \ptl_a \mathcal{E}^a \bigr) {\rm d}^2x{\rm d}t,
\end{align*}
where the phase-space is now
\[
 \{ (q, \mbo{\pi}), (\mathcal{A}_a, \mathcal{E}^a), (\gamma, p) \}
\]
with the Lagrange multipliers
\[
 \{ N, N^a, \mathcal{A}_0 \}
\]
and the constraints
\begin{subequations} \label{2+1constraints}
\begin{align}
&H= \bar{\mu}^{-1}_q \bigl(\Vert \mbo{\pi} \Vert^2_q - {\rm Tr}_q(\mbo{\pi})^2\bigr) + \frac{1}{8} p^2 + \halb {\rm e}^{-4\gamma}
q_{ab} \mathcal{E}^a \mathcal{E}^b + \bar{\mu}_q \bigl(-R_q + 2 q^{ab} \ptl_a \gamma \ptl_b \gamma\bigr) \notag\\
&\hphantom{H=}{}+ \frac{1}{4} {\rm e}^{4\gamma} q^{ab}q^{bd} \ptl_{[b} \mathcal{A}_{a]} \ptl_{[d} \mathcal{A}_{c]}, \\
&H_a= -2 \leftexp{(q)}{ \grad}_b \mbo{\pi}^b_a + p \ptl_a \gamma + \mathcal{E}^b \bigl(\ptl_{[a}\mathcal{A}_{b]} \bigr),\\
&\ptl_a \mathcal{E}^a =0.
\end{align}
\end{subequations}
It may be noted that $\mathcal{E}$ and $p$ are the momenta conjugate to $\mathcal{A} $ and $\gamma$ respectively.
After applying the Poincar\'e lemma on $\mathcal{E}$ and introducing the twist potential such that $\mathcal{E}^a =: \eps^{ab} \ptl_b \omega$,
we transform into the phase space
\[ X_{{\rm EWM}} = \bigl\{ (\gamma, p), (\omega, \mbo{r}), \bigl(q_{ab}, \mbo{\pi}^{ab}\bigr) \bigr\},
\]
where $\mbo{r}$ is the momentum conjugate to $\omega$. The variational principle reduces to
\begin{subequations} \label{2+1-no-constraints}
\begin{align}
&\tilde{J} := \int^{t_2}_{t_1}\int_{\Sigma} \bigl( \mbo{\pi}^{ab} \ptl_t q_{ab} + p \ptl_t \gamma + r \ptl_t \omega - N H - N^a H_a \bigr) {\rm d}^2x {\rm d}t,
\intertext{where $H$ and $H_a$ are now}
&H= \bar{\mu}^{-1}_q \biggl( \Vert \mbo{\pi} \Vert^2_q - {\rm Tr}_q (\mbo{\pi})^2 + \frac{1}{8} p^2 + \frac{1}{2} {\rm e}^{4\gamma} \mbo{r}^2 \biggr) \notag\\
&\hphantom{H=}{} + \bar{\mu}^{-1}_q \biggl(-R_q + 2 q^{ab} \ptl_a \gamma \ptl_b \gamma + \halb {\rm e}^{-4\gamma} q^{ab} \ptl_a \omega \ptl_b \omega \biggr), \\
&H_a= -2 \leftexp{(q)}{\grad}_b \mbo{\pi}^b_a + p \ptl_a \gamma + \mbo{r} \ptl_a \omega.
\end{align}
\end{subequations}
